\section{Detailed Setup Taxonomy}

\subsection{Setup Family Overview}

\begin{longtable}{p{0.24\linewidth}p{0.24\linewidth}p{0.44\linewidth}}
\toprule
\textbf{Setup family} & \textbf{Primary expression} & \textbf{Main edge hypothesis} \\
\midrule
Visible-liquidity bounce & Reversal from density / defended level & A real participant is willing to absorb opposing flow at a known price, allowing a tight stop and an asymmetric mean-reversion or rotation trade. \\
Visible-liquidity breakout & Consumption through density / level & Displayed defense is weaker than the market expected, and stop fuel plus empty space accelerate the move. \\
False breakout / failed continuation & Poke through level then return & Stops are harvested but no durable initiative flow remains, trapping late breakout traders. \\
Iceberg behavior & Hidden-liquidity defense or exhaustion & Refreshed hidden interest reveals a stronger participant than the visible book implies. \\
Absorption & Aggressive flow fails to move price materially & Initiative traders are being met by a larger passive participant. \\
Robot behavior & Repetitive algorithmic posting or lifting & The trader can temporarily piggyback on a detectable machine-driven process before it exhausts. \\
Spring / stretched-state reversal & One-way burst into anchor & A stretched local state snaps back once it meets meaningful resistance or defense. \\
News / volatility harvesting & Reaction trading after catalyst & News creates movement and opportunity; the trade is based on observed reaction, not semantic forecasting. \\
Leader-follower / guide instrument & Correlated or linked-instrument read-through & One instrument reveals the pressure or likely direction for another. \\
Cluster / PoC anchored logic & Executed-volume memory at price & Prior heavy business makes a zone behaviorally meaningful beyond current displayed size. \\
\bottomrule
\end{longtable}

\subsection{Universal Live Search Workflow}

Before specializing into individual setups, the materials imply a simple but disciplined live workflow:
\begin{enumerate}
  \item identify the next meaningful interaction price;
  \item classify the approach as drift, compression, impulse, repeated test, or stretched final push;
  \item watch the book and tape precisely at the interaction point;
  \item decide whether the event is becoming bounce, breakout, failure, absorption, or noise;
  \item enter only if stop, first target, and immediate post-entry behavior are all understandable.
\end{enumerate}

This matters because the setups below are not independent cartoon pictures.
Most of them are alternative interpretations of the same local confrontation at a level.

\subsection{Bounce From Large Visible Liquidity or Defended Density}

\setupfield{What the setup is.}{Price approaches a visible large resting order or battery of nearby resting orders and fails to move through it cleanly, producing a short-horizon reversal or bounce.}

\setupfield{Why it is believed to work.}{The trade assumes a real participant is defending a price, or at least that enough traders believe the defense is real that price reacts before the level is fully consumed. The attractiveness of the setup comes from a precise invalidation point: if the defense disappears or gets eaten, the idea is wrong quickly.}

\setupfield{Structural edge hypothesis.}{A defended level creates asymmetric risk because the trader can define the stop just beyond the defending participant, while the path away from the level may extend to the next cluster, opposing density, or session reference point.}

\setupfield{Best regime.}{Liquid instruments with enough volatility to move after the defense holds, but not such chaotic conditions that the level is constantly skipped over. The setup is especially favored in the source material when the anchor is clear and the stop is very short.}

\setupfield{How to search for it in real time.}{Start from a price that already matters before the touch: prior high/low, local extreme, cluster/PoC area, or obviously large visible size. Then watch \emph{how} price arrives. A cleaner bounce candidate usually forms after one-direction approach into the level or after a stretch/spring into the level, not in random chop. As the market touches the zone, ask four practical questions: Is the size still there? Is it being refreshed? Can aggressive traders move through it? Does price leave the area quickly once the interaction occurs?}

\setupfield{When to enter.}{Entry is typically near the defending level after evidence that the order is still functioning: repeated hits without collapse, supportive liquidity behind it, or a visible iceberg/refresh pattern. Some traders enter proactively into the level; others wait for a first rejection and then join the reversal.}

\setupfield{When not to enter.}{Do not enter just because a large number is visible. Avoid isolated densities in empty books, late entries after the bounce already traveled, or cases where there is no supportive liquidity behind the visible order and therefore no clean stop anchor.}

\setupfield{What to observe in the book.}{Size concentration at one or several nearby prices, persistence under attack, whether surrounding same-side liquidity remains supportive, whether opposite-side liquidity thins after rejection, and whether the visible level is static or keeps shifting.}

\setupfield{What to observe in tape and prints.}{Large aggressive prints hitting the level without clean penetration are supportive for the bounce thesis. If aggressive flow keeps landing and the level still holds, hidden or stronger-than-visible defense becomes more plausible.}

\setupfield{What to observe in price behavior.}{Repeated approach, hesitation, failure to gain acceptance beyond the level, and quick response away from the defended zone strengthen the case.}

\setupfield{Typical development.}{Price approaches, trades into the level, pauses or compresses, then rotates away. The best examples travel far enough to justify partial profit on a first target and leave runner potential toward the next meaningful obstacle.}

\setupfield{Typical confirming patterns.}{Alignment with chart support or resistance, nearby cluster PoC, iceberg refresh, spring-like prior stretch into the level, and strong opposing prints failing to make progress.}

\setupfield{Typical invalidation patterns.}{The defense vanishes, gets eaten quickly, or holds initially but price does not leave the area and instead grinds through. Another invalidation is a bounce that cannot even clear minor nearby obstacles, showing weak response quality.}

\setupfield{Likely traps and fake versions.}{A density can be spoofed, can be too small for the day regime, or can attract price only to disappear. Another trap is buying a ``bounce'' after price already bounced and the easy part is gone.}

\setupfield{Stop logic.}{Stop belongs just beyond the defended participant or the last structurally meaningful same-side support. If the trader needs a wide discretionary stop, the setup is no longer the clean short-stop bounce described in the source material.}

\setupfield{Exit and partial logic.}{Common practice is to take a first partial near the first reliable opposing obstacle and then decide whether the move has enough initiative to continue. The materials prefer taking what the market offers over demanding a heroic target.}

\setupfield{Failure modes.}{Visible defense without hidden support, wrong size calibration for the instrument, taking the trade against a stronger higher-level context, or entering after the level has already been weakened by too many tests.}

\setupfield{Clean versus unclean examples.}{A clean bounce has an obvious participant, a short stop, and visible room away from the level. An unclean bounce is based on vague hope, ambiguous book structure, or an anchor that is too small or too late.}

\setupfield{Explicit versus fuzzy parts.}{Explicit pieces include the visible anchor, stop location, and local reaction. Fuzzy pieces include how much hidden size is really there, how many tests are too many, and what counts as enough room for continuation.}

\subsection{Breakout Through Density or Through a Well-Defined Level}

\setupfield{What the setup is.}{Price pushes through a known level or dense visible order and then continues immediately into open space.}

\setupfield{Why it is believed to work.}{The source logic says breakouts work when the market has both a trigger and fuel. The trigger is the actual consumption of a visible barrier or obvious level. The fuel is a mix of stop orders, trapped traders, and absence of immediate opposing liquidity behind the level.}

\setupfield{Why the setup is considered harder than a bounce.}{The materials are blunt here: breakout trades are more difficult, faster, and less forgiving. A bounce gives a clearer stop. A breakout often requires instant judgment about whether the move is real.}

\setupfield{Structural edge hypothesis.}{A level with public attention accumulates conditional order flow. Once the level gives way, continuation can be fast if there is true initiative pressure and the book behind the level is thin.}

\setupfield{Best regime.}{Instruments that are actually moving, with clear larger-timeframe references and enough intraday energy to sustain expansion. The best examples occur when the market was already compressing into the level or when a visible defending order is visibly getting eaten.}

\setupfield{How to search for it in real time.}{First locate levels where a breakout would actually matter: range edge, session high/low, repeated intraday ceiling/floor, or a large visible density that many traders are clearly leaning on. Then watch whether the market approaches in a tightening and persistent way rather than in exhausted chaos. The trader is looking for signs that the barrier is becoming vulnerable: repeated probing, shallower pullbacks, same-side quotes stepping closer, and a lack of heavy fresh opposition behind the level.}

\setupfield{When to enter.}{Typical entries are just before or during the final consumption of the defending order, or on the first acceptance beyond the level if continuation is immediate. Delay is costly because a bad breakout should be recognized very fast.}

\setupfield{When not to enter.}{Do not trade breakouts in dead instruments, inside cluttered books, or into nearby opposing liquidity. Avoid levels that are obvious on the chart but unsupported by actual microstructure.}

\setupfield{What to watch in the book.}{Whether the opposing density is static or weakening, whether same-side support steps closer, whether there is empty space behind the level, and whether new opposing size appears immediately beyond the trigger point.}

\setupfield{What to watch in tape and prints.}{Initiative flow should hit the level with speed and continue after the breach. If the level is crossed but the tape loses urgency immediately, the breakout thesis weakens sharply.}

\setupfield{What to watch in price behavior.}{Compression, repeated testing, tight range into the level, then immediate post-break acceleration. A slow lazy drift above the level is less convincing than a fast acceptance.}

\setupfield{Typical confirming patterns.}{Repeated tests, shrinking pullbacks, visible last resistance getting consumed, correlated leaders moving the same way, and clear stop territory beyond the level.}

\setupfield{Typical invalidation patterns.}{No instant follow-through, fast return into the range, or visible aggressive buying/selling that fails to keep the market beyond the breached area.}

\setupfield{Likely traps and fake versions.}{Late participation after the move already launched, breakout into hidden opposing liquidity, and ``chart-only breakouts'' where the book never supported the idea.}

\setupfield{Stop logic.}{Because breakout trades can fail quickly, the stop usually belongs near the failed acceptance point, back inside the old range, or behind the final consumed barrier. If the move stalls instantly, exit may happen even before a hard stop is reached.}

\setupfield{Exit logic.}{Good breakouts should work soon. Partial exits can be taken into the first post-break opposing density or extension target. If continuation does not develop quickly, the source logic favors leaving rather than giving the trade too much room.}

\setupfield{Interaction with context.}{Breakouts are stronger when the day already rewards trend extension, when leaders align, when news expands movement, or when a prior trapped side is clearly visible.}

\setupfield{Algorithmically important ambiguity.}{The human notion of ``it must go immediately'' is central and should eventually become a measurable post-trigger acceptance window.}

\subsection{False Breakout and Failed Continuation}

\setupfield{What the setup is.}{Price breaches a visible level or range boundary, triggers breakout participation or stop orders, and then returns back through the level instead of continuing.}

\setupfield{Why it is believed to work.}{The false breakout setup exploits the asymmetry between public attention and true initiative flow. The crowd sees the trigger. The market briefly uses that trigger for liquidity. But if no durable demand or supply exists beyond it, the move reverses and traps the late participants.}

\setupfield{Core logic.}{A breakout level can fail because the visible barrier was not the real obstacle, because the stop run exhausted the move, because a larger participant used the excursion to complete inventory, or because broader context did not support continuation.}

\setupfield{How to search for it in real time.}{Do not search for false breakouts everywhere. Search only at levels where many people are likely to expect continuation. Then, once the level breaks, ignore the picture and focus on the post-break behavior: does the market accept price on the new side, or does it look uncomfortable there? The discretionary trader is effectively watching for a short sequence: breach, weak continuation, loss of urgency, return, trap.}

\setupfield{What to observe.}{Fast poke through a level, poor follow-through, aggressive flow that does not carry price further, and a decisive return into the prior range.}

\setupfield{Supporting clues.}{Repeated failed acceptance above/below the breached area, cluster history showing heavy business but no clean expansion, or a prior spring that makes the move look more like exhaustion than healthy initiation.}

\setupfield{When to enter and when to skip.}{Entry is usually after the market clearly re-enters the prior zone or fails multiple times to gain acceptance beyond the break. Skip the trade if the move is simply retesting in an orderly way, if a strong leader instrument still supports continuation, or if the tape remains decisively one-sided beyond the level.}

\setupfield{How humans recognize cleaner examples.}{Cleaner false breakouts often feel ``too easy'' to breakout traders: the level is obvious, the public trigger is visible, but the market cannot sustain price on the other side.}

\setupfield{Stop logic.}{Stop generally belongs beyond the excursion extreme. Once the market re-enters the prior range convincingly, the invalidation point becomes much tighter than in a blind countertrend fade.}

\setupfield{Failure mode.}{The main danger is mistaking a normal shallow retest for failure. Entering against a true breakout too early is one of the fastest ways to get run over.}

\setupfield{What remains fuzzy.}{The difference between brief pause and genuine failure depends on pace, tape quality, and day regime. That timing threshold is a major research target.}

\subsection{Iceberg-Like Behavior}

\setupfield{What the setup is.}{A price level shows repeated replenishment or volume-throughput inconsistent with the visible resting amount, suggesting hidden liquidity.}

\setupfield{Why it is believed to work.}{Icebergs reveal that the visible book understates true interest. If the trader correctly identifies a real iceberg, the level may be defended much more strongly than naive observers think.}

\setupfield{Observed markers.}{Repeated refill of visible size, cluster or traded volume far exceeding the displayed amount, and persistence at one price despite continued interaction.}

\setupfield{How to search for it in real time.}{The trader is usually not scanning for icebergs in empty space. Search begins at a suspiciously important level where prints keep landing but the visible quote does not behave normally. The practical clue is throughput mismatch: too much business for too little displayed size. On a good screen the trader keeps asking whether the level should already have disappeared if the visible size were all that existed.}

\setupfield{How it is traded.}{Two broad ways appear in the materials:
\begin{itemize}
  \item trade \emph{from} the iceberg as a defended anchor;
  \item trade the \emph{failure} or exhaustion of the iceberg when it finally gives way.
\end{itemize}}

\setupfield{Important nuance.}{The sources repeatedly warn that iceberg-highlighting software is helpful but imperfect. A highlight is not a proof. Visual confirmation and contextual understanding remain necessary.}

\setupfield{When not to trust it.}{Possible false positives include fast reposting bots, fragmented visible quotes, and situations where throughput is large simply because the market is extremely active.}

\setupfield{Stop and exit logic.}{If leaning on an iceberg, the stop belongs beyond the hidden participant's presumed defense zone. If trading its failure, the stop belongs back on the defended side after the market proves the iceberg is gone.}

\setupfield{What makes a clean iceberg trade.}{The cleaner version is not just an indicator highlight but a whole story: meaningful level, repeated refresh, large real interaction, and a visible reaction once the iceberg either holds or finally fails.}

\setupfield{Research difficulty.}{Iceberg detection is especially hard in order-book-only data, because the strongest evidence comes from the relationship between displayed size and executed volume.}

\subsection{Liquidity Absorption}

\setupfield{What the setup is.}{Aggressive orders continue to hit one side, but price does not travel proportionally. The market appears to be absorbing initiative flow.}

\setupfield{Why it matters.}{Absorption often precedes either reversal or delayed breakout, depending on who is actually winning. It is therefore not a directional setup by itself but a critical interpretation layer.}

\setupfield{Bullish or bearish?}{Absorption against aggressive buyers can be bearish if they cannot lift price. Absorption against aggressive sellers can be bullish if they cannot press price lower. But if the absorbing side is merely a staging area before continuation, the same picture can become a launching pad.}

\setupfield{How to search for it in real time.}{Search for absorption only near a level that matters. Then compare pace of aggressive prints to actual displacement. If a lot of market buying hits offers but price barely rises, or heavy selling hits bids but price barely falls, the trader begins to suspect absorption. The next decision is not direction yet; it is whether the passive side is truly in control or merely delaying the inevitable.}

\setupfield{What to watch.}{Tape speed versus price progress, repeated hits at one price, whether the quote refreshes, and whether the absorbing zone eventually starts to repel price or instead suddenly caves.}

\setupfield{How it usually develops.}{Price stalls near a key zone while one side keeps pressing. Then either:
\begin{itemize}
  \item the aggressor exhausts and price rotates away; or
  \item the passive side is finally overwhelmed and the delayed break becomes forceful.
\end{itemize}}

\setupfield{Why this setup is hard.}{Absorption is observationally rich but interpretation-heavy. The exact same tape can look like strong defense or like pre-break build-up. Context decides.}

\subsection{Tape-Based Aggression and Impulse Confirmation}

\setupfield{What the setup is.}{Aggressive tape itself is not always the trigger, but it serves as confirmation that a local microstructure event has real initiative behind it.}

\setupfield{Use in practice.}{The tape is used to answer questions such as:
\begin{itemize}
  \item is the level being attacked seriously or only lightly;
  \item are prints speeding up;
  \item is the market lifting offers or hitting bids in meaningful size;
  \item is a move supported by real aggression or only by quotes shifting around.
\end{itemize}}

\setupfield{How to search for it in real time.}{Tape aggression is not usually hunted as a standalone setup at first glance. The trader first finds a location of interest, then uses the tape to decide whether the developing bounce or breakout deserves trust. In practice the tape is a truth-serum layer: is this move being powered by actual market orders or only by displayed quotes shuffling around?}

\setupfield{How it combines with other setups.}{Tape aggression is a confirmation layer for breakouts, a stress test for defended levels, and a warning sign when a robot or visible density is about to stop working.}

\setupfield{Key trap.}{Aggression without context can simply be the final emotional burst into a larger passive participant. Tape is a \emph{force} measure, not an automatic direction measure.}

\subsection{Robot or Algorithm Behavior in the Book}

\setupfield{What the setup is.}{A repetitive algorithmic pattern becomes visible through posting cadence, equal-size updates, predictable stepping, or repeated passive or aggressive behavior.}

\setupfield{Why it is believed to work.}{The trader does not need to know the robot's ultimate objective. The exploitable part is that the robot temporarily makes behavior more regular and therefore more forecastable.}

\setupfield{Typical signs.}{Plus-one stepping, repetitive same-size refreshing, continuous passive accumulation, or machine-like pushing that looks too regular for discretionary behavior.}

\setupfield{How to search for it in real time.}{The trader watches for unnatural regularity: equal-size reposting, same cadence, same step size, or a quote that ``walks'' price in a mechanical way. The setup becomes tradable only if the machine behavior is currently dominating the local book enough that other traders begin reacting to it.}

\setupfield{How it is traded.}{The material strongly suggests ``trade with the robot while it is still working.'' That means piggybacking while the machine is still pressing or supporting, then exiting as soon as its influence weakens, gets absorbed, or the move becomes too stretched.}

\setupfield{When not to trust it.}{Late recognition, weak robot footprint, crowded already-extended move, or obvious counterpressure. Robots evolve; there is no eternal static taxonomy.}

\setupfield{Failure mode.}{A trader falls in love with the robot idea, holds after the behavior changes, and forgets that the edge existed only while the machine's pattern remained active.}

\subsection{POC / Cluster-Anchored Setups}

\setupfield{What the setup is.}{Executed-volume memory, especially PoC alignment, is used as an additional anchor for bounce or continuation logic.}

\setupfield{Why it is believed to work.}{A zone with heavy executed business can mark genuine prior interest rather than merely displayed intent. When current book structure aligns with historical heavy business, the level gains credibility.}

\setupfield{Strong version.}{Current density aligns with chart structure \emph{and} recent cluster concentration. The same zone matters to both present and past participation.}

\setupfield{Weak version.}{Current density appears but clusters show little historical business there, making the apparent level less trustworthy.}

\setupfield{How to search for it in real time.}{Before the session or during quiet moments, mark where heavy prior business sits. Then, when live liquidity appears near those prices, ask whether present displayed intent and past executed business are pointing to the same zone. Traders use this to avoid overreacting to single fresh quotes that have no historical memory behind them.}

\setupfield{Research implication.}{Cluster-derived features should eventually complement live order-book and tape features, especially for Russian instruments where footprint-style memory was emphasized in the teaching material.}

\subsection{Spring or Stretched-State Reversal}

\setupfield{What the setup is.}{A one-direction move extends so far without normal pullback that even a modest real counterforce can produce a sharp snapback.}

\setupfield{Why it is believed to work.}{Once a move is stretched, participants are vulnerable: late chasers enter at bad prices, earlier participants may start taking profit, and a defending participant needs less help to trigger rotation.}

\setupfield{Critical nuance.}{The spring is not the reason by itself. The materials explicitly warn that ``the spring alone'' is not sufficient. The trader still needs a real anchor such as density, iceberg, or meaningful structural level.}

\setupfield{How to search for it in real time.}{Look for a move that has become one-directional, emotionally extended, and increasingly expensive to chase. Then ask whether it is now running into something concrete: visible defense, prior session level, cluster memory, or leader divergence. Without that concrete anchor, ``it feels stretched'' is treated as a bad reason to trade.}

\setupfield{Best examples.}{Fast extension into a clear opposing zone, visible local exhaustion, then quick reaction.}

\setupfield{Main trap.}{Trying to fade every strong move just because it feels stretched. Many strong trends keep going much farther than expected.}

\subsection{News and Volatility-Harvesting Trades}

\setupfield{What the setup is.}{A news event or known catalyst creates a temporary expansion in movement and opportunity. The trade is based on the actual reaction in the book and tape, not on semantic prediction of the news headline.}

\setupfield{Why it is believed to work.}{Catalysts increase participation, expand range, and create temporary disorder in which both breakout and reversal structures can become unusually profitable.}

\setupfield{Practical rule from the material.}{Do not try to be a macro commentator. Treat news as a reason the market may move and then trade the observable microstructure.}

\setupfield{How to search for it in real time.}{When a catalyst is scheduled or obvious, do not guess direction first. Instead search for the first clean interaction between volatility and structure: fast break through a known level, violent overshoot into real liquidity, or clear one-sided tape that actually survives beyond the first impulse. The materials are much closer to ``trade the reaction'' than ``predict the headline.''}

\setupfield{Best uses.}{Breakout confirmation, fast impulse participation, volatility harvesting when the book becomes temporarily one-sided, or selective bounce trades after a news-driven overshoot into real liquidity.}

\setupfield{Main dangers.}{Slippage, skipped levels, fake liquidity, and emotionally attractive but low-information chaos. News is opportunity, but also a fast path to undisciplined trading.}

\subsection{Leader-Follower, Guide-Instrument, and Correlation Logic}

\setupfield{What the setup is.}{One instrument or market acts as a guide for another. The source materials mention this in currencies, futures, and broader correlated structures.}

\setupfield{Why it is believed to work.}{Some instruments react first, or express the relevant pressure more cleanly. If the traded instrument is noisier or thinner, the leader can help interpret whether a local move is trustworthy.}

\setupfield{Examples from the materials.}{Currency and futures complexes, RTS versus currency relationships, oil as a partial leader, and cases where a local DOM cannot be trusted in isolation because external references dominate the instrument.}

\setupfield{What to watch.}{Short-lag directional agreement, divergence, whether the follower is delayed, and whether the local book behavior is consistent with what the leader suggests.}

\setupfield{How to search for it in real time.}{The trader keeps a small set of trusted leaders on screen and uses them selectively. A leader is most useful when the traded instrument becomes locally ambiguous. If the local DOM is messy but the guide instrument is pressing a level cleanly, that can decide whether to trust a breakout, reject a fade, or wait.}

\setupfield{Research difficulty.}{This requires synchronized multi-instrument data and careful timestamp alignment. It is promising, but more demanding than single-instrument order-book studies.}

\subsection{Special Regime: Limit-Up / Limit-Down and Extreme Exchange Constraints}

\setupfield{What the setup is.}{Exchange limits create a highly distorted environment in which normal microstructure logic becomes unreliable or dangerous.}

\setupfield{Why it matters.}{The source material treats these situations as special-risk regimes rather than normal pattern territory. Traders can become trapped with limited exit routes, one-sided books, and violent squeezes.}

\setupfield{Practical lesson.}{Do not apply ordinary discretionary templates mechanically inside exceptional exchange-constraint regimes. Special handling and conservative participation are required.}
